\frontsection{Introduction Générale}

Dans le cadre de ma formation en cycle ingénieur, ce stage de fin d'année porte sur la conception et le développement d'une application web dédiée à la gestion des abonnés et des projets d'une association. L'objectif est de proposer une solution centralisée permettant de faciliter le suivi des adhérents, des abonnements, des projets associatifs et des différentes opérations financières qui leur sont associées.

\vspace{0.5cm}

Le fonctionnement de l'association repose notamment sur les cotisations annuelles de ses abonnés ainsi que sur différentes sources de financement destinées à la réalisation de projets à caractère social, tels que le forage de puits ou la construction et la rénovation de bâtiments. Pour chaque projet, un comité est constitué afin d'assurer son suivi de sa création jusqu'à sa clôture. Les membres de ce comité doivent pouvoir enregistrer les dons, les dépenses et les différentes pièces justificatives liées aux opérations réalisées.

\vspace{0.5cm}

La gestion de ces informations nécessite une organisation rigoureuse et une bonne traçabilité des opérations financières. L'application développée doit ainsi permettre de centraliser ces données, de gérer les droits d'accès des différents membres, de conserver les justificatifs associés aux opérations et de faciliter la préparation des rapports des projets. Elle doit également permettre à chaque abonné de consulter son propre historique, notamment ses abonnements, ses dons et les éventuelles sommes restant dues.

\vspace{0.5cm}

La problématique de ce stage consiste donc à concevoir une solution adaptée aux besoins de l'association, permettant d'améliorer la gestion, le suivi et la traçabilité de ses activités tout en simplifiant l'accès aux informations pour les différents acteurs.

\vspace{0.5cm}

Ce rapport présente les différentes étapes de réalisation de ce projet. Le premier chapitre présente l'organisme d'accueil et le contexte du stage. Le deuxième chapitre est consacré à l'étude et à l'analyse des besoins. Enfin, le troisième chapitre présente la conception et la mise en œuvre de la solution développée.
